\section{Methods}
\label{sec:methods}

\subsection{Population Dataset}
\label{sec:dataset}
We analyzed a longitudinal dataset of coagulation, fibrinolysis, thrombin generation
assay (TGA), and rotational thromboelastometry (ROTEM) measurements from women
desiring a pregnancy. The study was approved by the Institutional Review Board at
the University of Vermont. All participants gave written consent before enrollment.
Blood collection, plasma processing, and assay protocols have been described
previously~\citep{McLean2012TGA,Hale2012,Psoinos2021,Bernstein2016}. Briefly, citrated platelet-poor plasma
was collected without tourniquet use after supine rest at three clinical visits:
Visit~1 (pre-pregnancy baseline in the follicular phase), Visit~2 (end of first
trimester), and Visit~3 (mid-third trimester). Plasma aliquots were stored at
$-80^{\circ}$C for subsequent analysis. Measurements of secondary analytes were
measured as follows: Stago clotting assays measured coagulation-factor activity, ELISA
measured fibrinolysis proteins, and the CLIA-certified laboratory at the University of Vermont
Medical Center measured hormone levels. We performed TGA using 5~pM relipidated tissue factor with
fluorogenic substrate on a Synergy4 plate reader as described
in McLean et al.\citep{McLean2012TGA}. Clot formation and fibrinolysis dynamics were assessed via
viscoelastometry using the ROTEM Delta Instrument (Werfen). Thawed plasma was mixed
with or without 4~nM tPA (with respect to plasma volume), recalcified with 15~mM
calcium chloride, and incubated for 3~minutes at 37$^{\circ}$C. The reaction was
initiated by adding the plasma mixture to ROTEM cups containing tissue factor at a
17:1 ratio; the final concentration of TF in the reaction was 5~pM.

The dataset contained 153 total records across approximately 50 subjects. After
removing columns with $>30\%$ missing values and retaining only subjects with
complete data across all three visits, we obtained $K=23$ complete longitudinal
profiles across $n=72$ assay measurements per visit.
The 72 assay measurements spanned five categories: (i)~hormones (estradiol,
progesterone), (ii)~coagulation factors and inhibitors (Factors~II, V, VII, VIII,
IX, X, XI, XII; antithrombin (AT), protein~C (PC), tissue factor pathway inhibitor (TFPI),
von Willebrand factor (vWF), fibrinogen, etc.), (iii)~fibrinolytic markers
(plasminogen, plasminogen activator inhibitor-1 (PAI-1), plasminogen activator inhibitor-2 (PAI-2),
$\alpha_2$-antiplasmin ($\alpha_2$-AP), thrombin-activatable fibrinolysis inhibitor (TAFI),
tissue plasminogen activator (tPA), D-dimer equivalents),
(iv)~thrombin generation parameters under four initiator conditions (TF at 5~pM,
TF+thrombomodulin (TM), No~TF, No~TF+anti-FXIa), and (v)~ROTEM/viscoelastic parameters
(clotting time (CT), maximum clot firmness (MCF),
alpha angle, lysis onset time (LOT), maximum lysis (ML), lysis time (LT),
area under the curve (AUC) under TF~Only and TF+tPA conditions).

Demographics and clinical characteristics of the resulting cohort, including
age at enrollment, prepregnancy BMI, parity, gestational ages at each visit,
and the cross-tabulation of enrollment cohort against pregnancy outcome, are
summarized in Table~\ref{tab:demographics}. The 23 patients were drawn from
three enrollment conditions: Healthy nulliparous women ($n{=}14$), women with
a personal history of preeclampsia from a previous pregnancy (Prior~PE,
$n{=}6$), and women with polycystic ovarian syndrome (PCOS, $n{=}3$;
2~individuals were nulliparous). For conditional generation, we defined three
overlapping subgroups based on pregnancy outcome and comorbidity rather than
enrollment condition: Uncomplicated ($n{=}18$, all patients whose study pregnancy
did not result in preeclampsia), PCOS ($n{=}3$, a cross-cutting comorbidity;
1~patient also developed PE), and Developed~PE ($n{=}5$, patients who developed
preeclampsia during the study pregnancy, drawn from 2~healthy-enrolled,
2~prior-PE-enrolled, and 1~PCOS-enrolled participants).


\subsection{Multiplicity-Weighted Stochastic Attention}
\label{sec:sa}

We generated synthetic patients using a multiplicity-weighted variant of Stochastic
Attention (SA) rooted in modern Hopfield network theory~\citep{Ramsauer2021}. We
stored $K$ memory patterns as columns of a matrix $\mhat{X}\in\R^{d\times K}$ and
defined a weighted Hopfield energy landscape:
\begin{equation}
    E_r(\vxi) = \frac{1}{2}\norm{\vxi}^2
    - \frac{1}{\beta}\log\sum_{k=1}^{K} r_k \exp\left(\beta\,\mathbf{m}_k^\top\vxi\right)
    \label{eq:hopfield_energy}
\end{equation}
where $\{\mathbf{m}_k\}_{k=1}^K$ are stored memory patterns, $\beta$ is an inverse
temperature parameter controlling the sharpness of retrieval, and $\{r_k\}_{k=1}^K$
are per-pattern multiplicity weights. When all $r_k = 1$, this reduces to the standard
modern Hopfield energy; when $r_k = \rho > 1$ for a designated subset of patterns, the
energy landscape is continuously deformed to favor that subset.

\rone{This energy defines a probability distribution that can be written exactly. Writing
$\pi_r(\vxi)\propto\exp[-\beta E_r(\vxi)]$ and completing the square in each exponent gives
the expression:
\begin{equation}
    \exp\!\left[-\frac{\beta}{2}\norm{\vxi}^2 + \beta\,\mathbf{m}_k^\top\vxi\right]
    = \exp\!\left(\frac{\beta}{2}\norm{\mathbf{m}_k}^2\right)
      \exp\!\left[-\frac{\beta}{2}\norm{\vxi-\mathbf{m}_k}^2\right],
    \label{eq:complete_square}
\end{equation}
Substituting Eq.~\eqref{eq:complete_square} into the probability density gives an exact
finite isotropic Gaussian mixture:
\begin{equation}
    \pi_r(\vxi) = \sum_{k=1}^{K} q_k\,
    \mathcal{N}\!\left(\mathbf{m}_k,\ \beta^{-1}\mathbf{I}\right),
    \qquad
    q_k \propto r_k\exp\!\left(\frac{\beta}{2}\norm{\mathbf{m}_k}^2\right).
    \label{eq:exact_mixture}
\end{equation}
Each stored patient is the center of one component. The common component variance is
$\beta^{-1}$, and the multiplicity weights enter the mixture probabilities directly. The
memories used here are normalized to unit length, so the exponential factor is shared by
every component and $q_k = r_k/\sum_j r_j$ exactly. We could therefore sample this
distribution directly: first choose $k$ from $q$ and then draw
$\vxi\sim\mathcal{N}(\mathbf{m}_k,\beta^{-1}\mathbf{I})$. This procedure requires no Markov
chain. The cohort
analyzed in this work was generated with the Langevin implementation given next, and that is
what the reported results describe. We added the direct sampler for the stochastic-attention
target afterward as a check,
using the same normalization, magnitude rescaling, and decoding steps
(Table~\ref{stab:sampling-ladder}).}

We sampled from this landscape using an Unadjusted Langevin Algorithm (ULA):
\begin{equation}
    \vxi_{t+1} = (1-\alpha)\,\vxi_t
    + \alpha\,\mhat{X}\,\text{softmax}\!\left(\beta\,\mhat{X}^\top\vxi_t + \log\mathbf{r}\right)
    + \sqrt{\frac{2\alpha}{\beta}}\,\boldsymbol{\varepsilon}_t
    \label{eq:ula}
\end{equation}
where $\alpha$ is a step size, $\boldsymbol{\varepsilon}_t\sim\mathcal{N}(\mathbf{0},\mathbf{I})$,
and $\log\mathbf{r}$ is the element-wise log of the multiplicity vector. The
deterministic component drove $\vxi_t$ toward a multiplicity-weighted combination of
stored patterns, while the stochastic component added variation. The
$\log r_k$ bias in the softmax logits interpolated between unconditioned generation ($\rho = 1$,
all patterns equally weighted) and hard subset curation ($\rho \to \infty$, only
designated patterns contribute). This is an inference-time operation that requires
no change to the stored patterns.

\rone{The synthetic cohort analyzed in this work was produced by initializing each Langevin
chain near a randomly selected stored pattern and retaining the final state after $T{=}2000$
steps. Because the starting point and chain length could affect the output, we compared four
sampling procedures. These used the original initialization near a stored pattern, random
initialization on the unit sphere, pooled samples from multiple chains after burn-in and
thinning, and Metropolis-adjusted sampling. As a separate check, we also drew directly from
the exact stochastic-attention mixture in Eq.~\eqref{eq:exact_mixture}, without using a Markov
chain. The fidelity and
membership-inference results were similar across all four procedures and the direct draws.
This agreement indicated that the original implementation reproduced the evaluated behavior
of the sampling distribution at $\beta^*$ (Supplementary
Fig.~\ref{sfig:sampling-ladder}).}

\rboth{The inverse temperature $\beta$ affects Hopfield retrieval and direct sampling in
different ways. In the Hopfield update, small $\beta$ spreads attention across many stored
profiles and tends toward their weighted average, whereas large $\beta$ concentrates attention on
individual profiles. When sampling directly from Eq.~\eqref{eq:exact_mixture}, the stored
profile is chosen according to $q_k$. For the unit-normalized profiles used here, these
probabilities are independent of $\beta$ and are set by the multiplicity weights. Instead,
$\beta$ controls the spread around the selected profile through the covariance
$\beta^{-1}\mathbf{I}$. Small $\beta$ therefore produces broader draws, whereas large $\beta$
produces tighter draws. The attention-entropy transition used to define $\beta^*$ marks the
shift from attention shared across profiles to attention concentrated on individual profiles.}
\rboth{We selected $\beta^*$ from the transition in normalized attention entropy across a
logarithmic sweep of $\beta$. Specifically, $\beta^*$ was the grid point where the entropy
curve bent downward most strongly, calculated as the most negative finite-difference second
derivative with respect to $\log\beta$ (Supplementary Methods). This choice depended on the
relative positions of the $K$ profiles in $d$ dimensions, not on what their features
represented. For this dataset
($K{=}23$, $d_{\text{PCA}}{=}18$), it gave $\beta^*\approx2.94$; the random-pattern prediction
$\sqrt{d}\approx4.24$ served as a theoretical reference. A sensitivity analysis across
$\beta/\beta^*\in[0.1,3.0]$ showed that the selected value balanced novelty and
fidelity (Table~\ref{stab:beta-sweep}). For conditioned generation, we repeated the entropy
calculation with the multiplicity bias included, giving a $\rho$-dependent $\beta^*(\rho)$.}

\rboth{We used the participation ratio to measure how strongly multiplicity weighting
concentrated repeated draws around a few stored profiles. For unit-normalized memories,
$q_k=r_k/\sum_j r_j$ and
$K_{\text{eff}}=1/\sum_k q_k^2=(\sum_k r_k)^2/\sum_k r_k^2$. Thus,
$K_{\text{eff}}$ is the number of equally weighted components with the same concentration.
Equal weights gave $K_{\text{eff}}=23$; concentrating weight on three profiles moved it toward
three. This was not an independent clinical sample size. To assign a target fraction
$f_{\text{target}}$ of component probability to a subgroup, we used
$\rho=f_{\text{target}}K_{\text{bg}}/[K_{\text{des}}(1-f_{\text{target}})]$, where
$K_{\text{des}}$ and $K_{\text{bg}}$ are the subgroup and background profile counts.}


\subsection{Pipeline for Continuous Longitudinal Data}
\label{sec:pipeline}

Applying SA to continuous longitudinal clinical data required several adaptations
beyond the standard Hopfield sampling framework. We first concatenated each patient's
$n{=}72$ assay measurements across all three visits into a single vector of dimension
$d_{\text{concat}} = 3n = 216$, so that each stored memory pattern encoded a
complete longitudinal profile and generated synthetic patients would have internally
consistent multi-visit trajectories. We then standardized each feature to zero mean
and unit variance and applied PCA retaining 95\% of the variance, reducing
dimensionality from 216 to $d_{\text{PCA}}{=}18$. This compression served two
purposes: it removed collinearity among the 216 features and yielded a
memory-to-dimension ratio of $K/d_{\text{PCA}} \approx 1.28$, placing SA in a
favorable operating regime where the number of stored patterns exceeded the
dimensionality of the representation.

Standard SA operates on unit-norm patterns, which is appropriate for discrete data
(e.g., one-hot protein sequences) but destroys the anisotropic variance structure
of continuous clinical measurements, collapsing dispersion across all principal
components to a unit sphere. We addressed this with a direction-magnitude
decomposition. Before normalization, we recorded each pattern's PCA-space norm
$s_k = \|\mathbf{m}_k\|$; we then ran SA on unit-norm patterns to obtain a
directional sample $\hat{\vxi}$, drew a magnitude $s$ from the empirical
distribution of $\{s_k\}$, and reconstructed the rescaled sample as
$\vxi_{\text{rescaled}} = s \cdot \hat{\vxi}$. This preserved the directional
structure learned by the Hopfield energy while restoring the natural scale of
variation. The rescaled PCA vector was mapped back to the original feature space via
inverse PCA and destandardization, then split into three 72-dimensional visit records.

To generate condition-specific cohorts (e.g., PCOS-only), we set the multiplicity
$r_k = \rho$ for patterns belonging to the target subgroup and $r_k = 1$ for all
others, and sampled using the weighted Langevin update (Eq.~\ref{eq:ula}). The
multiplicity $\rho$ was computed to place a target fraction
$f_{\text{target}} = 0.80$ of the mixture-component probability on the designated subset,
while retaining 20\% probability on background patterns, via
$\rho = f_{\text{target}} \cdot K_{\text{bg}} / (K_{\text{des}} \cdot (1 -
f_{\text{target}}))$. Multiplicity parameters for all three subgroups are
listed in Table~\ref{stab:multiplicity}. For the PCOS subgroup
($K_{\text{des}}{=}3$), this required $\rho \approx 26.7$, reducing the
participation ratio to $K_{\text{eff}} \approx 4.6$. \rone{For both unconditioned and
conditioned generation, magnitudes were drawn from the empirical distribution of all $23$
PCA-space norms. Subgroup weighting affected directional sampling, not the distribution of
magnitudes.}


The complete generation pipeline is summarized in Algorithm~\ref{alg:sa-pipeline}.
The full set of SA hyperparameters used for all experiments is listed in
Table~\ref{stab:hyperparams}, with key operating values $\alpha = 0.01$,
$T = 2{,}000$ Langevin iterations per sample, and random seed 42 for
reproducibility. A control experiment comparing ULA with the
Metropolis-Adjusted Langevin Algorithm (MALA) showed that the
discretization bias is negligible at this step size. Across 10 chains of
5{,}000 iterations the MALA acceptance rate was $99.9 \pm 0.03\%$, and mean
post-burn-in energies and effective sample sizes were indistinguishable
between the two samplers (Supplementary Table~\ref{tab:ula-vs-mala}). The pipeline was implemented in Julia (v1.12) using the packages listed in the code
repository. Generating 100 synthetic patients required approximately 0.1~seconds on
a standard laptop (Apple M-series), comparable to MVN sampling ($\approx$0.15~s),
because SA operates in the 18-dimensional PCA space rather than the full
216-dimensional feature space. Each synthetic patient was generated by an independent
Langevin chain (no shared state between patients); generation quality was stable
across $T \in [500, 4000]$ iterations (median MRE 1.0--1.6\%), indicating that
$T{=}2{,}000$ was sufficient for convergence.

\begin{algorithm}[ht]
\caption{Multiplicity-Weighted SA for Longitudinal Patient Generation}
\label{alg:sa-pipeline}
\begin{algorithmic}[1]
\Require Patient data $\{\mathbf{x}_k^{(v)}\}$ for $k=1,\ldots,K$ patients, $v=1,2,3$ visits
\Require Multiplicity vector $\mathbf{r}$ (all ones for unconditioned; $r_k = \rho$ for designated subset)
\Ensure $N$ synthetic longitudinal patient profiles
\State Concatenate: $\mathbf{x}_k \gets [\mathbf{x}_k^{(1)}; \mathbf{x}_k^{(2)}; \mathbf{x}_k^{(3)}] \in \mathbb{R}^{216}$
\State Standardize: $\mathbf{z}_k \gets (\mathbf{x}_k - \boldsymbol{\mu}) \oslash \boldsymbol{\sigma}$
\State PCA: $\mathbf{m}_k \gets \mathbf{W}^\top \mathbf{z}_k \in \mathbb{R}^{18}$ \Comment{95\% variance retained}
\State Record norms: $s_k \gets \|\mathbf{m}_k\|$; normalize $\hat{\mathbf{m}}_k \gets \mathbf{m}_k / s_k$
\State \rboth{Select $\beta^*$ where the attention-entropy curve bends downward most strongly, including $\log\mathbf{r}$ when conditioned}
\For{$i = 1, \ldots, N$}
    \State \rone{Select anchor $k_0$ from all patterns, or from the designated subset when conditioned}
    \State \rone{$\hat{\vxi}_0 \gets \hat{\mathbf{m}}_{k_0}+0.01\boldsymbol{\varepsilon}_0$; normalize $\hat{\vxi}_0$}
    \For{$t = 0, \ldots, T-1$} \Comment{Langevin dynamics}
        \State $\hat{\vxi}_{t+1} \gets (1-\alpha)\hat{\vxi}_t + \alpha\hat{\mathbf{M}}\,\mathrm{softmax}(\beta^*\hat{\mathbf{M}}^\top\hat{\vxi}_t + \log\mathbf{r}) + \sqrt{2\alpha/\beta^*}\,\boldsymbol{\varepsilon}_t$
    \EndFor
    \State Draw magnitude: $s \sim \{s_k\}_{k=1}^K$
    \State Rescale: $\mathbf{m}_{\text{synth}} \gets s \cdot \hat{\vxi}_T / \|\hat{\vxi}_T\|$
    \State Inverse PCA + destandardize: $\mathbf{x}_{\text{synth}} \gets \boldsymbol{\sigma} \odot (\mathbf{W}\,\mathbf{m}_{\text{synth}} + \bar{\mathbf{z}}) + \boldsymbol{\mu}$
    \State Split into visit records: $\mathbf{x}_{\text{synth}}^{(1)}, \mathbf{x}_{\text{synth}}^{(2)}, \mathbf{x}_{\text{synth}}^{(3)}$
\EndFor
\end{algorithmic}
\end{algorithm}

\subsection{Baseline and Validation Framework}
\label{sec:validation}

As a baseline, we generated synthetic patients by fitting a single multivariate
normal (MVN) distribution to the same $K{=}23$ concatenated 216-dimensional patient
profiles. Both methods operated on the same concatenated
representation. Each patient's three visits were joined into one vector before
generation, so that both methods had the opportunity to capture cross-visit
dependencies. For SA, this structure was preserved through the Hopfield energy
landscape operating on 18-dimensional PCA projections of the concatenated vectors.
For MVN, the $216 \times 216$ sample covariance matrix had rank at most~22 and was
regularized using Ledoit--Wolf shrinkage~\citep{Ledoit2004}, which inflated
eigenvalues in the 194 null dimensions.
\rone{To examine the role of the shared PCA representation, we also evaluated a fitted
two-component diagonal-covariance Gaussian mixture, a kernel-density estimator,
$k$-nearest-neighbor interpolation, and a Gaussian copula in the same 18-dimensional PCA
space. Their implementations and shared evaluation pipeline are described in Supplementary
Methods (PCA-space ablation).}
We then evaluated both SA and MVN synthetic patients through four progressively stringent
validation levels.
\textbf{Level~1 (Marginal Plausibility)} tested whether individual feature
distributions matched, using per-feature means, standard deviations, and known
physiological relationships. \textbf{Level~2 (Joint Structure)} tested whether the
cross-visit covariance was preserved, by comparing full $216 \times 216$ correlation
matrices, eigenvalue spectra, and PCA projections between real, SA, and MVN
populations. \textbf{Level~3 (Conditional Structure)} tested whether rare subgroups
could be amplified while preserving condition-specific signatures, by generating
cohorts of 100 patients each from the Uncomplicated ($n{=}18$), PCOS ($n{=}3$), and
Developed~PE ($n{=}5$) subgroups using multiplicity-weighted sampling. \textbf{Level~4 (Mechanistic
Consistency)} tested whether synthetic patients produced biologically plausible
outputs under \rboth{a separately specified, generator-blind mechanistic model calibrated
on the same real cohort}. We used the Hockin--Mann BZ2012
coagulation model~\citep{Hockin2002,BrummelZiedins2012}, a system of 58 ODEs with 64 rate
constants, in which 5 rate constants were calibrated on Visit~1 real patients at the
population level and the remaining 59 were held at literature values. We ran the model
on all real and synthetic patients under two conditions (TF-only and TF+TM) and
compared predicted-versus-measured thrombin generation features. The primary metric
was cloud overlap: the fraction of synthetic predicted/measured ratios falling
within the 5th--95th percentile range of real ratios.

\rone{For the membership-inference analysis, we drew $40$ random partitions of the $23$ real
profiles. For each partition, we repeated the preprocessing using $15$ profiles, constructed
the memory matrix, and generated $100$ synthetic profiles; the other $8$ real profiles were
held out. We assigned each real profile a membership score equal to the negative Euclidean
distance to its nearest synthetic profile in standardized $216$-dimensional space, so a
higher score indicated stronger evidence that the profile had been stored. A receiver
operating characteristic curve compared the fraction of stored profiles correctly identified
with the fraction of held-out profiles incorrectly identified as the score threshold changed.
We calculated the area under this curve using the equivalent Mann--Whitney rank statistic,
with ties contributing one half, and reported the mean across the $40$ partitions. An AUC of
$0.5$ indicated chance discrimination, whereas $1.0$ indicated perfect discrimination.}

\subsection{\rboth{Robustness Tests}}
\label{sec:robustness-benchmarks}

\rboth{Results from one clinical cohort could not show how stochastic-attention generation
behaved as training-set size, dimension, and nonlinear structure changed. We therefore used
three tests. The covariance test compared stochastic attention with a sample-covariance
Gaussian generator on low-rank Gaussian populations with known covariance. We generated one
training cohort for every tested combination from the factor model
$\mathbf{X}=\mathbf{Z}\mathbf{L}^{\top}+\boldsymbol{\epsilon}$, where
$\mathbf{Z}\in\mathbb{R}^{n\times r}$ and $\mathbf{L}\in\mathbb{R}^{p\times r}$ had
independent standard-normal entries and $\boldsymbol{\epsilon}$ had independent Gaussian
entries with variance $0.1$. This gave the known population covariance
$\boldsymbol{\Sigma}_{\mathrm{pop}}=\mathbf{L}\mathbf{L}^{\top}+0.1\mathbf{I}$. We tested
every combination of
$n\in\{8,15,23,40,80\}$, $p\in\{30,120,216\}$, and
$r\in\{3,5,10\}$. For each training cohort, we standardized the features, retained enough
principal components to explain $95\%$ of the variance, constructed the stochastic-attention
memory matrix from the resulting scores, and selected $\beta$ separately using the same
attention-entropy criterion as in the clinical analysis. Stochastic attention generated
$100$ profiles using ULA with $\alpha{=}0.01$, $T{=}2{,}000$, and empirical magnitude
rescaling. The Gaussian generator used the training mean and sample covariance. For numerical
sampling, diagonal jitter began at $10^{-8}\mathbf{I}$ and increased tenfold as needed through
$10^{-2}\mathbf{I}$; if Cholesky sampling still failed, we sampled in the nonnegative
eigenspace of the sample covariance. We applied no covariance shrinkage. The Gaussian
generator also produced $100$ profiles. Covariance error was
$\lVert\widehat{\boldsymbol{\Sigma}}_{\mathrm{gen}}-
\boldsymbol{\Sigma}_{\mathrm{pop}}\rVert_F/
\lVert\boldsymbol{\Sigma}_{\mathrm{pop}}\rVert_F$. Figure~\ref{fig:sim-recovery}A reports
Gaussian error divided by stochastic-attention error, averaged across the three latent ranks.
A ratio above one favored stochastic attention.}

\rboth{The remaining tests asked how stochastic attention changed with less clinical
information and how it compared with the Gaussian generator on a nonlinear population. For
the sample-size test, we drew ten subsets without replacement from the $23$ complete patients
at each $n\in\{20,15,10,8\}$. For every subset, we repeated the standardization, PCA, and
$\beta$ selection, constructed a new stochastic-attention memory matrix, generated $100$
profiles, and evaluated the generated cohort against the same $23$-patient reference. The
outcomes were relative cross-visit correlation error, pooled median marginal relative error,
and TF-only mechanistic cloud overlap. Figure~\ref{fig:sim-recovery}B shows the mean and
standard deviation of cross-visit error across the ten subsets. For the nonlinearity test, we
drew $t$ uniformly from $[-1,1]$ and defined an S-curve by $x_1=t$ and
$x_2=c\sin(\pi t)$. Viewed in the $(x_1,x_2)$ plane, this one-dimensional path resembled the
letter S because it had two opposing bends. Setting $c{=}0$ produced a straight line, and
increasing $c$ made the bends stronger. Each simulated profile was a noisy point near this
path. For each replicate, we used a new random transformation to spread the two coordinates
across $p{=}30$ or~$120$ simulated variables while preserving the curve's geometry. We then
added Gaussian noise with standard deviation $0.05$ in every direction. We used $n{=}23$ and curvature
$c\in\{0,0.25,0.5,1,1.5,2\}$, with five replicates per setting. From every training cohort, we
constructed a stochastic-attention memory matrix and estimated the same sample-covariance
Gaussian comparator; each generator produced $100$ profiles. We projected generated profiles
onto the known two-dimensional embedding plane and measured their mean distance from the
curve, thereby excluding noise orthogonal to that plane. We also measured relative covariance
error against the covariance of an independent $50{,}000$-profile population draw. For
the reported curvature threshold, we averaged stochastic-attention curve distance across the
five replicates and both dimensions at each curvature and identified the first nonzero value at
which this mean was at least twice its value for the straight-line population. All simulation
cells used deterministic, cell-specific random-number streams derived from seed~42.}
